% !TeX root = ../navier-stokes-manuscript.tex
\subsection{The spatial row inside the first response}\label{sub:TheViscosityDrivenSpatialAscent}

The table has placed temporal response against spatial height.  Start with the
base response \(V_0=u\), where the two directions first became visible at
critical height.  For \(s\geq0\), set
\[
  E_s(t):=\frac12\norm{\Lambda^su(t)}_2^2
\]
and retain the complete transport--pressure pairing
\[
  \mathcal P_s(t)
  :=-
  \operatorname{Re}
  \pair{\Lambda^su}
  {\Lambda^s\bigl((u\cdot\nabla)u+\nabla p\bigr)}_{L^2}.
\]
The spatial energy row is
\begin{equation}\label{eq:BodyFullSpatialEnergyBalance}
  E_s'=\mathcal P_s-2\nu E_{s+1}.
\end{equation}
At \(s=\tfrac12\),
\begin{equation}\label{eq:BodyFirstCriticalRow}
  H'=\mathcal P_{1/2}-2\nu E_{3/2}.
\end{equation}
The next two differentiated balances make the triangular structure explicit:
\begin{align}
  H''
  &=\partial_t\mathcal P_{1/2}
  -2\nu\mathcal P_{3/2}
  +(2\nu)^2E_{5/2},
  \label{eq:BodySecondCriticalRow}\\
  H'''
  &=\partial_t^2\mathcal P_{1/2}
  -2\nu\partial_t\mathcal P_{3/2}
  +(2\nu)^2\mathcal P_{5/2}
  -(2\nu)^3E_{7/2}.
  \label{eq:BodyThirdCriticalRow}
\end{align}
Every lower transfer survives while one higher spatial energy appears.

The complete identity keeps every lower transfer while exposing one new
spatial height.
\begin{proposition}[Critical-height spatial ascent]\label[proposition]{prop:CriticalHeightSpatialAscent}
For every integer \(n\geq1\),
\begin{equation}\label{eq:BodyCriticalHeightSpatialAscent}
  H^{(n)}
  =(-2\nu)^nE_{n+1/2}
  +\sum_{j=0}^{n-1}
  (-2\nu)^j
  \partial_t^{\,n-1-j}\mathcal P_{j+1/2},
  \qquad n\geq1.
\end{equation}
\end{proposition}
\begin{proof}
The critical balance \eqref{eq:BodyFirstCriticalRow} is the case \(n=1\).
Assume the identity at order \(n\), differentiate it, and apply
\eqref{eq:BodyFullSpatialEnergyBalance} at height \(n+\tfrac12\).  The new
transfer extends the sum by one term, while the viscous contribution becomes
  \((-2\nu)^{n+1}E_{n+3/2}\).  This is the identity at order \(n+1\).
\end{proof}
Rearranging for the highest energy gives
\begin{equation}\label{eq:BodyCompletedCriticalRow}
  E_{n+1/2}
  =\frac{
    H^{(n)}-
    \displaystyle\sum_{j=0}^{n-1}
    (-2\nu)^j\partial_t^{\,n-1-j}\mathcal P_{j+1/2}
  }{(-2\nu)^n}.
\end{equation}
The division is lawful because \(\nu>0\).  It does not erase the transfer
sum.  This identity names the exact collection that has to be controlled
before the exposed spatial energy is controlled.

The height derivative remains a contraction of the velocity responses:
\begin{equation}\label{eq:BodyHeightDerivativeVelocityContraction}
  H^{(n)}
  =\frac12\sum_{a=0}^{n}\binom{n}{a}
  \operatorname{Re}
  \pair{\Lambda^{1/2}V_a}
       {\Lambda^{1/2}V_{n-a}}_{L^2}.
\end{equation}
The three readings at order \(n\) are distinct.  \(V_n\) is the Eulerian
response, \(H^{(n)}\) is its quadratic critical contraction, and
\(E_{n+1/2}\) is the highest spatial energy exposed after the full transfer
chain has been retained.  The identity reveals that energy; it does not bound
the transfer chain on its right.  This is the exact point at which the hoped-for
reversal becomes conditional.  Viscosity has exposed a higher spatial energy;
writing the identity has not controlled that energy.

Does the same obstruction remain after the base response changes?  Shift the
calculation from \(u\) to a fixed \(V_r\).  The Laplacian still acts linearly on
that response, while every split of the transport derivative remains in the
binomial sum.  The same calculation now answers the question for all finite
\(r\).
